\chapter{מודלים מאומנים מראש}

\section{הרעיון של הדרכה עצמית}

מודלים מאומנים מראש (\textenglish{pre-trained models}) הם מודלים שכבר נלמדו על כמות ענקית של טקסט.
ההדרכה הזו נעשית בדרך כלל על משימה "עצמית", כמו חיזוי האסימון הבא במשפט או השלמה של מילים חסרות.

\section{\textenglish{BERT} וההדרכה הדו־כיוונית}

\textenglish{BERT} (Bidirectional Encoder Representations from \textenglish{Transformers})
הוא מודל מאומן מראש שמצביע לעבר דו־כיווני, כלומר הוא רואה את כל ההקשר.
זה הוא בעלת יתרון על מודלים חד־כיווניים בעבור משימות כמו סיווג טקסט וחילוץ מידע.

\section{\textenglish{GPT} וההדרכה החד־כיוונית}

\textenglish{GPT} (\textenglish{Generative Pre-trained Transformer})
מאומנת בדרך אחת בלבד: חיזוי האסימון הבא.
למרות זה, היא מסוגלת ליצור טקסט קוהרנטי בעל אורך רב.

\section{ההדרכה בתנאי תרמוספרה}

כאשר משתמשים במודל מאומן מראש ליישום ספציפי חדש, אנו משתמשים בתהליך בשם \textenglish{fine-tuning}.

\begin{equation}
\text{Loss}_{fine-tune} = \lambda \cdot \text{Loss}_{task} + (1 - \lambda) \cdot \text{Loss}_{pretrain}
\end{equation}

כאן \(\lambda\) הוא היפרפרמטר המסדיר את האיזון בין השניים.

\section{גדולים יותר טוב?}

מחקרים רבים הראו שמודלים גדולים יותר מציגים "פריצות" בביצועים,
שבהן הם פתאום מצליחים בתרגילים שהם כשלו בהם קודם.
זה נקרא \textenglish{emergent abilities} של דוגמנות שפה גדולה.

\section{דוגמאות למודלים מאומנים מראש}

כמה מודלים מאומנים מראש פופולריים כוללים:

\begin{enumerate}
\item \textenglish{BERT} – דו־כיווני, טוב לסיווג
\item \textenglish{GPT} – חד־כיווני, טוב ליצירה
\item \textenglish{RoBERTa} – שיפור של \textenglish{BERT}
\item \textenglish{LLaMA} – מודל אחרי של \textenglish{Meta}
\item \textenglish{T5} – מודל סימטרי יחסית
\end{enumerate}

\section{יישום של מודל קיים}

אתה יכול להשתמש בספריה \textenglish{Hugging Face} להטעין מודל קיים:

\begin{LTR}
\begin{verbatim}
from transformers import AutoTokenizer, AutoModel

tokenizer = AutoTokenizer.from_pretrained("bert-base-hebrew")
model = AutoModel.from_pretrained("bert-base-hebrew")

inputs = tokenizer("שלום עולם", return_tensors="pt")
outputs = model(**inputs)
\end{verbatim}
\end{LTR}

קוד זה מטעין את \textenglish{BERT} לעברית ומעבד טקסט עברי כלשהו.
